\chapter{Plan d’expérience}

\section{Objectif du plan}
Construire une campagne expérimentale reproductible pour identifier les facteurs dominants.

\section{Facteurs et niveaux (version initiale)}
\begin{longtable}{@{}llll@{}}
\toprule
Facteur & Symbole & Niveaux initiaux & Remarque \\
\midrule
\endhead
Pression d’alimentation & $P_{in}$ & bas / moyen / haut & Contrôle source \\
Ouverture d’anche & $h_0$ & 3 niveaux & Réglage mécanique \\
Longueur de conduit & $L$ & 2 à 3 niveaux & Effet résonateur \\
Température & $T$ & contrôlée & Condition d’essai \\
\bottomrule
\end{longtable}

\section{Variables de réponse}
\begin{itemize}
  \item Fréquence fondamentale et harmoniques.
  \item Niveau de pression acoustique (SPL).
  \item Seuil d’amorçage et extinction des oscillations.
  \item Stabilité temporelle (jitter, dérive fréquentielle).
\end{itemize}

\section{Plan recommandé pour démarrer}
\begin{enumerate}
  \item Criblage factoriel fractionnaire pour identifier les facteurs dominants.
  \item Plan de surface de réponse sur les facteurs principaux.
  \item Répétitions au point central pour estimer la variabilité.
\end{enumerate}
